\originalchapter{8.701 作业五:中微子与核稳定}
题源: Fall 2020 Problem Set 5, 2020年11月4日. 本章三题均有官方答案; 原题 1(d) 的角度表述须加明确更正.
\begin{exercise}\label{ex:8701-5-1}
原题 1. T2K 利用 $\pi^+\to\mu^+\nu_\mu$ 的离轴束流. Pion 沿实验室 z 轴以 $\beta$ 运动, 中微子在 pion 静止系能量 $E^*$, 极角 $\theta^*$. (a) 求 $E^*=(m_\pi^2-m_\mu^2)/(2m_\pi)$. (b) 证明 $E=\gamma E^*(1+\beta\cos\theta^*)$ 及 $E\cos\theta=\gamma E^*(\cos\theta^*+\beta)$. (c) 证明 $\gamma^2(1-\beta\cos\theta)(1+\beta\cos\theta^*)=1$. (d) 原题要求证明``实验室最大角 $\theta_{\max}=1/\gamma$''. (e) 在 $\theta\ll1$ 时证明 $E\simeq0.43E_\pi/(1+\beta\gamma^2\theta^2)$. (f) Pion 能谱在 1--5 GeV 平坦, 画 $2.5^\circ$ 离轴的中微子能谱, 说明距束流 295 km 的 Super-Kamiokande 为何取此角.
\end{exercise}
\begin{solution}
(a)--(c) 官方框架重编. 两体运动学给 $E^*=29.8\,\mathrm{MeV}$, boost 后能量和纵向动量即 (b). 无质量中微子 $p_z=E\cos\theta$, 消去 $\theta^*$ 即得 (c), 同时 $E=E^*/[\gamma(1-\beta\cos\theta)]$.

(d) 编者更正. 该命题作为全局最大角不成立: $\theta^*=\pi$ 时 $\cos\theta=-1$, 实验室角可到 $\pi$. $1/\gamma$ 是前向特征角; 对固定离轴角的能量极大值, 高 boost 近似满足 $\gamma\theta\simeq1$. 不提供对错误命题的伪证明.

(e) $1-\beta\cos\theta\simeq(1-\beta)+\beta\theta^2/2$, $1-\beta\simeq1/(2\gamma^2)$, 因而所求成立, $0.43$ 来自 $1-m_\mu^2/m_\pi^2\simeq0.427$. 这是高 boost 及小角联合近似.

(f) 令 $a=\theta^2/m_\pi^2$, $k=0.427$, 则 $E_\nu=kE_\pi/(1+aE_\pi^2)$, 最大值为 $k/(2\sqrt a)\simeq0.683\,\mathrm{GeV}$, 对应 $E_\pi\simeq3.20\,\mathrm{GeV}$. 1、5 GeV 两端分别约 $0.389,0.620\,\mathrm{GeV}$. 转换后的数密度需乘雅可比 $\sum|\dd E_\pi/\dd E_\nu|$ 及角发射权重 $1/[4\pi\gamma^2(1-\beta\cos\theta)^2]$. 在极大值附近多个 pion 能量映到相近中微子能量, 形成窄峰; 理想零角宽模型有平方根堆积, 有限束流角宽与能量分辨率使它平滑. 不能把平坦 pion 谱直接画成平坦中微子谱.
\begin{center}\begin{tikzpicture}
\begin{axis}[width=8cm,height=4.8cm,xlabel={$E_\pi$ (GeV)},ylabel={$E_\nu$ (GeV)},domain=1:5,samples=120,ymin=0.3,ymax=0.75]
\addplot[thick,structurecolor]{0.427*x/(1+(0.0436332/0.13957)^2*x^2)};
\end{axis}\end{tikzpicture}\end{center}
上图为能量映射, 下图补出实际的理想分箱能谱.
\begin{center}\begin{tikzpicture}\begin{axis}[width=9cm,height=5cm,xlabel={$E_\nu$ (GeV)},ylabel={归一分箱密度 (GeV$^{-1}$)},xmin=0.38,xmax=0.70,ymin=0]
\addplot[thick,structurecolor,const plot] coordinates {(0.3840,0.030457) (0.3880,0.299385) (0.3920,0.308500) (0.3960,0.317811) (0.4000,0.325254) (0.4040,0.335932) (0.4080,0.343648) (0.4120,0.354723) (0.4160,0.363826) (0.4200,0.374203) (0.4240,0.384779) (0.4280,0.395555) (0.4320,0.405326) (0.4360,0.417681) (0.4400,0.429054) (0.4440,0.441895) (0.4480,0.453709) (0.4520,0.465721) (0.4560,0.479269) (0.4600,0.493082) (0.4640,0.507161) (0.4680,0.520094) (0.4720,0.536088) (0.4760,0.550959) (0.4800,0.567574) (0.4840,0.583017) (0.4880,0.600250) (0.4920,0.617816) (0.4960,0.635712) (0.5000,0.653937) (0.5040,0.674121) (0.5080,0.694708) (0.5120,0.714008) (0.5160,0.737049) (0.5200,0.758825) (0.5240,0.782755) (0.5280,0.807151) (0.5320,0.832009) (0.5360,0.859170) (0.5400,0.888737) (0.5440,0.915114) (0.5480,0.947715) (0.5520,0.979041) (0.5560,1.010958) (0.5600,1.047482) (0.5640,1.086789) (0.5680,1.122772) (0.5720,1.165686) (0.5760,1.209477) (0.5800,1.256291) (0.5840,1.306225) (0.5880,1.359380) (0.5920,1.413602) (0.5960,1.475702) (0.6000,1.541392) (0.6040,1.610761) (0.6080,1.686261) (0.6120,1.770491) (0.6160,1.858903) (0.6200,6.409795) (0.6240,6.905012) (0.6280,7.149299) (0.6320,7.423562) (0.6360,7.736551) (0.6400,8.089562) (0.6440,8.499740) (0.6480,8.990459) (0.6520,9.569872) (0.6560,10.293025) (0.6600,11.209261) (0.6640,12.440331) (0.6680,14.201791) (0.6720,17.055863) (0.6760,23.030545) (0.6800,45.290342) (0.6840,0)};
\end{axis}\end{tikzpicture}\end{center}
这是平坦 1--5 GeV pion 谱、静止系各向同性发射、固定 $2.5^\circ$ 探测角的理想模型, 用完整 boost 和发射角权重做 4 MeV 分箱; 不是实测谱. 峰顶随分箱细化变尖, 不代表有限分辨率实验中会有无限计数.
对应谱在 $0.39$--$0.68\,\mathrm{GeV}$ 有支持、在 $0.683$ GeV 附近集中. $L/E$ 对应大气质量差首个振荡峰约 $0.57$ GeV, 与这个窄谱相近, 因而离轴能谱有利于振荡测量与抑制高能背景.
\end{solution}
\begin{exercise}\label{ex:8701-5-2}
原题 2. 用指定原子质量模型
\begin{equation*}
 M=(A-Z)m_n+Z(m_p+m_e)-a_vA+a_sA^{2/3}+a_c\frac{Z^2}{A^{1/3}}+a_a\frac{(A-2Z)^2}{4A}+\frac\delta{\sqrt A}
\end{equation*}
求固定 A 的最稳定 Z 并画 $Z(A)$, 解释各项. 取 $a_v=15.67,a_s=17.23,a_c=0.714,a_a=93.15$, 偶偶、奇 A、奇奇分别 $\delta=-11.2,0,+11.2$, 单位 MeV/$c^2$.
\end{exercise}
\begin{solution}
官方方法重编. 在每个配对分支内先连续求导, 配对项不依 Z. 记 $\Delta=m_n-m_p-m_e\simeq0.782\,\mathrm{MeV}$, 得
$Z_*=A(a_a+\Delta)/[2a_a+2a_cA^{2/3}]$.
$A=20,100,200$ 时连续值约 $9.54,43.27,79.90$. \begin{center}\begin{tikzpicture}\begin{axis}[width=8cm,height=4.5cm,xlabel={$A$},ylabel={$Z_*$},domain=2:250,samples=150,ymin=0]
\addplot[thick,structurecolor]{x*(93.15+0.782)/(2*93.15+2*0.714*x^(2/3))};
\end{axis}\end{tikzpicture}\end{center}
此图使用本题指定系数而非第 \ref{ch:nuclear-mass} 章示意系数. 真实整数核素须比较邻近两三项并加偶奇配对. 体积、表面、库仑、不对称和配对解释与式 \eqref{eq:semf} 相同, 但本题 $a_a$ 比常见版本大四倍, 不得漏掉分母 $4A$.
\end{solution}
\begin{exercise}\label{ex:8701-5-3}
原题 3. 天然铀含 $^{238}$U $99.28\%$、$^{235}$U $0.72\%$. 若形成时二者相等, $\tau_{235}=10^9$ 年, $\tau_{238}=6.6\times10^9$ 年, 物质年龄是多少? 自地壳形成 $2.5\times10^9$ 年以来, $^{238}$U 衰变了多少比例?
\end{exercise}
\begin{solution}
官方解答重编. 比值为 $N_{235}/N_{238}=\exp[-t(1/\tau_{235}-1/\tau_{238})]=0.72/99.28$. 解得 $t=\log(99.28/0.72)/(1/10^9-1/(6.6\times10^9))\simeq5.81\times10^9$ 年. 第二问 $1-\e^{-2.5/6.6}\simeq31.5\%$. 年龄依``初始数量相等、此后无混入或流失''的题设, 不把此值当作今日太阳系年龄的实测定论.
\end{solution}
